% Updated evidence-bounded conclusion, 7 November 2026.
\section{Research Summary and Answers}

This thesis develops B-SMART as an evidence-aware SME analytics and decision-support framework, supported by an operational multi-tenant software prototype and offline empirical experiments. The supplied retail data contain 32 shops labelled grocery/FMCG, not a verified cross-sector SME sample. Merchant origin, lawful collection and participant permissions have not been independently authenticated. The results below are therefore preliminary evidence for the supplied retail setting, not a population-wide claim.

\paragraph{RQ1: Demand forecasting.}
The original daily shop--category LightGBM L2 model achieved mean WAPE 81.12\% across three rolling-origin folds. A history-rich LightGBM L1 model reduced mean WAPE to 74.46\%, an improvement of 6.65 percentage points. However, the L1 model underpredicted total volume by approximately 32.5\% across the combined test predictions, so its lower WAPE alone does not make it suitable for inventory order quantities. The data are sparse and lumpy: 37.60\% of daily shop--category observations have zero recorded sales, and SKU-level daily series are even sparser.

The strongest operational improvement is to predict cumulative demand directly over the next 7, 14 and 28 days. Across three non-overlapping rolling-origin evaluation folds, the Tweedie models achieved WAPE of 33.62\%, 25.19\% and 19.08\%, respectively, compared with four-block historical baselines of 38.62\%, 28.90\% and 21.91\%. The corresponding Tweedie volume biases were -0.78\%, -2.55\% and +0.06\%. A four-fold held-out-shop test achieved 33.43\% 7-day WAPE when historical sales from the held-out shops were available but shop ID was withheld. This supports transfer within the same grocery/FMCG setting, not zero-shot cold start or transfer to other sectors. Category-specific errors remain higher for sparse categories such as Baby Care and Eggs. Quantile intervals under-covered their nominal 80\% target and require further calibration.

The defensible RQ1 conclusion is therefore not that daily SKU forecasting has been solved. Rather, the results show that the forecast decision horizon and aggregation level materially affect performance, and that direct lead-time-aligned demand forecasts are more useful candidates for inventory planning than a biased one-day point forecast. The original branch--SKU decision objective remains only partially evaluated because the current reliable daily panel is shop--category and SKU assortment/stock availability history is incomplete.

\paragraph{RQ2: Customer inactivity.}
The offline LightGBM inactivity model achieved PR-AUC 0.7405, compared with 0.7346 for logistic regression and 0.6528 for a recency rule. Its incremental advantage over logistic regression is modest. Aggregate TreeSHAP contributions identify historical purchase frequency, tenure and recency as influential model features. These associations do not establish causal retention benefit, incremental campaign response, or owner understanding.

\paragraph{RQ3: Strategy optimisation.}
The 2,500 supplied policy cases remain synthetic-oracle scenarios, not a replay using the new horizon-specific forecasts and verified historical inventory. The trend policy has lower simulated cost by BDT 19,400.60 but increases lost units from 45 to 100. As 97.68\% of cases begin with stock at least equal to assumed future demand, the scenario set has limited ability to distinguish policies. A real forecast-driven inventory evaluation still requires stock ledgers, receipts, supplier lead times, costs, and observed shortage or waste outcomes.

\paragraph{RQ4: Explainability and usability.}
The 1,800 structured explanation cases show that recommendation records can carry explicit reasons and constraints. The supplied evidence-complete flag is not independently verified. No owner-participant study has measured task correctness, comprehension, completion time or trust. Therefore technical explanation structure has been demonstrated, but human usability has not.

\section{Contribution}

The strongest evidence-supported contribution is an evidence-aware, multi-horizon decision-support design that connects validated business records, demand estimates at decision-relevant horizons, explicit operating constraints, explanation, owner choice and later outcome monitoring. Empirically, the study demonstrates that (i) daily shop--category forecasting is strongly affected by zero-demand and series heterogeneity; (ii) absolute-error optimisation can reduce WAPE while producing serious volume underprediction; (iii) direct 7-, 14- and 28-day cumulative-demand models substantially outperform same-horizon moving-average baselines; and (iv) held-out-shop transfer is feasible within the represented grocery/FMCG setting when historical sales are available. These findings are more defensible than a claim that one universal model is best for every SME.

\section{Remaining Work Before Journal Submission}

First obtain independent data-owner authorisation and document collection, consent and ethics requirements. Next reconcile assortment and stock-availability history so that zero sales can be distinguished from censored demand. Validate the selected horizon-specific model on a genuinely future or independently collected test period, and calibrate uncertainty intervals to the intended service level. Connect the forecast to RQ3 using real opening stock, incoming purchase orders, supplier lead time, pack-size constraints, holding/shortage costs and observed outcomes. Complete SKU allocation only after validating assortment history and SKU-level mix. Run consented SME-owner tasks for RQ4 and measure comprehension, decision correctness, completion time and trust. Deploy versioned models only after API integration, monitoring, drift checks, access control and rollback tests. Finally, test additional independent sectors before making cross-sector SME claims.

The thesis should preserve the negative findings, uncertainty and data boundaries. Lower WAPE at a longer cumulative horizon is not the same target as daily SKU accuracy, and a synthetic policy result is not realised business savings. A journal manuscript should focus on the evidence-supported multi-horizon forecasting contribution and its limitations rather than presenting the entire software platform as already field-validated.
